\definecolor{routingcolor}{RGB}{42,157,143}
\definecolor{monolithiccolor}{RGB}{231,111,81}
\newcommand{\metricbar}[2]{%
  \makebox[4.0cm][l]{\textcolor{black!12}{\rule{4.0cm}{9pt}}}%
  \hspace{-4.0cm}%
  \textcolor{#1}{\rule{#2}{9pt}}%
}

\chapter{EXPERIMENTAL STUDY}

This chapter documents the current state of the experimental study in two
complementary ways. First, it reports two completed MMLU experiments that
already provide a concrete architecture comparison under a shared setup.
Second, it clarifies which parts of the broader thesis evaluation pipeline are
already usable and which still remain preliminary. This separation is important
because the repository can already produce meaningful pilot evidence, but not
every benchmark surface is yet mature enough for thesis-grade claims.

\section{Evaluation Environment}

The platform uses a split deployment model in which the control plane runs locally while model inference runs on remote GPU hosts. Table \ref{tab:environment-matrix} summarizes the current environment used by the project.

\begin{table}[htbp]
    \centering
    \caption{Current experimental environment and its role in the platform.}
    \label{tab:environment-matrix}
    \small
    \begin{tabularx}{\linewidth}{@{} l l X @{}}
        \toprule
        \textbf{Layer} & \textbf{Technology} & \textbf{Role in the experiment stack} \\
        \midrule
        Local control plane & FastAPI, Next.js, MLflow, Python runner & Experiment launch, tracking, result storage, and UI management. \\
        Dedicated SLM serving & Dockerized vLLM on Google Cloud L4 machines & Low-cost SLM inference for baseline, hybrid, and swarm components. \\
        Shared mid-tier serving & Dockerized vLLM on a Google Cloud RTX PRO6000 host & Shared 20B--40B serving for escalation and stronger baselines. \\
        Heavy serving & Dockerized vLLM on Nebius H100/H200-class machines & Frontier-scale serving for the largest baselines and fallback paths. \\
        Measurement layer & Metrics proxy, resource profiles, MLflow, JSON/Markdown reporter & Latency, token, cost, energy, carbon, and experiment summary collection. \\
        \bottomrule
    \end{tabularx}
\end{table}

All selected models are exposed through OpenAI-compatible HTTP endpoints so that the experiment runner can stay provider-neutral. The repository also contains host-lock logic to coordinate experiments that depend on shared hosts where only one large model may be active at a time.

\section{Current MMLU Comparison Setup}

This report includes two new MMLU experiments and uses them to compare two
response-generation strategies under the same benchmark. The routing run,
\texttt{exp\_ecdd81b2}, first answers with the lightweight
\texttt{gemma4-4b} model and escalates uncertain cases to
\texttt{gemma4-31b}. The monolithic run, \texttt{exp\_018ab0cf}, answers every
question directly with \texttt{gemma4-31b}. Since both runs evaluate the same
\texttt{mmlu} benchmark with 100 samples and zero temperature, the comparison
isolates the effect of architecture choice rather than dataset variation.

The comparison focuses on five outcome families: accuracy, EATS, latency,
cost, and resource usage. In this repository, EATS is the
``Efficiency-Accuracy Trade-off Score'' and rewards architectures that preserve
accuracy while reducing heavy LLM usage and normalized cost. Therefore, a
higher EATS indicates a better efficiency-quality balance.

\begin{table}[htbp]
\centering
\small
\renewcommand{\arraystretch}{1.2}
\caption{Configuration of the two MMLU comparison runs}
\label{tab:mmlu-config}
\begin{tabular}{|p{3.3cm}|p{5.1cm}|p{5.1cm}|}
\hline
\textbf{Setting} & \textbf{Routing: \texttt{exp\_ecdd81b2}} & \textbf{Monolithic: \texttt{exp\_018ab0cf}} \\
\hline
Architecture & Routing with SLM draft and LLM fallback & Single-pass monolithic inference \\
\hline
Benchmark & \multicolumn{2}{c|}{MMLU} \\
\hline
Samples & \multicolumn{2}{c|}{100} \\
\hline
Final LLM & \texttt{gemma4-31b} fallback only & \texttt{gemma4-31b} for every sample \\
\hline
Draft SLM & \texttt{gemma4-4b} & None \\
\hline
Temperature & \multicolumn{2}{c|}{0.0} \\
\hline
Confidence threshold & 0.8, active for escalation & 0.7 recorded in config, not used for routing \\
\hline
LLM call ratio & 34\% & 100\% \\
\hline
\end{tabular}
\end{table}

\section{Score Comparison}

Figure~\ref{fig:mmlu-scores} compares the two architectures on raw task
quality. The monolithic design achieved the higher MMLU accuracy with 85
correct answers out of 100, while the routing design reached 76 correct
answers. This 9 percentage point gap is expected because the monolithic
baseline always uses the stronger \texttt{gemma4-31b} model.

The more interesting result is the EATS comparison. Routing reached an EATS of
2.171, while the monolithic baseline remained at 0.842. This means the routing
architecture delivered a 2.58 times better efficiency-accuracy trade-off even
though its raw accuracy was lower. The main driver is selective escalation:
66 samples were accepted at the SLM stage and only 34 required the heavy LLM.

\begin{figure}[htbp]
\centering
\begin{minipage}[t]{0.46\textwidth}
\centering
\textbf{Accuracy (\%, higher is better)}\\[0.6em]
\begin{tabular}{@{}l l r@{}}
Routing & \metricbar{routingcolor}{3.58cm} & 76.0 \\
Monolithic & \metricbar{monolithiccolor}{4.00cm} & 85.0 \\
\end{tabular}
\end{minipage}
\hfill
\begin{minipage}[t]{0.46\textwidth}
\centering
\textbf{EATS (higher is better)}\\[0.6em]
\begin{tabular}{@{}l l r@{}}
Routing & \metricbar{routingcolor}{4.00cm} & 2.171 \\
Monolithic & \metricbar{monolithiccolor}{1.55cm} & 0.842 \\
\end{tabular}
\end{minipage}
\caption{Score comparison between routing and monolithic architectures on MMLU}
\label{fig:mmlu-scores}
\end{figure}

\begin{table}[htbp]
\centering
\small
\renewcommand{\arraystretch}{1.2}
\caption{Raw outcome summary for the two MMLU runs}
\label{tab:mmlu-summary}
\begin{tabular}{|l|c|c|}
\hline
\textbf{Metric} & \textbf{Routing} & \textbf{Monolithic} \\
\hline
Accuracy & 76.0\% & 85.0\% \\
\hline
EATS & 2.171 & 0.842 \\
\hline
Correct answers & 76 / 100 & 85 / 100 \\
\hline
LLM call ratio & 34\% & 100\% \\
\hline
Average latency & 20.45 s & 9.29 s \\
\hline
Latency p95 & 42.18 s & 17.66 s \\
\hline
Total cost & \$0.5913 & \$0.6066 \\
\hline
Total energy & 0.0665 kWh & 0.0826 kWh \\
\hline
Total CO$_2$ & 25.28 g & 31.39 g \\
\hline
Total tokens & 69,512 & 42,889 \\
\hline
\end{tabular}
\end{table}

\section{Efficiency and Resource Comparison}

Figure~\ref{fig:mmlu-efficiency} shows the operating-cost side of the same
comparison. Routing reduced heavy-model usage from 100\% to 34\%, lowered total
energy from 0.0826 to 0.0665 kWh, and reduced estimated CO$_2$ emissions from
31.39 g to 25.28 g. The cost difference was smaller but still favorable to
routing: \$0.5913 versus \$0.6066, or roughly a 2.5\% reduction.

The latency trade-off is much sharper. Average latency increased from 9.29
seconds in the monolithic run to 20.45 seconds in routing, and tail latency
(p95) increased from 17.66 seconds to 42.18 seconds. This overhead comes from
the two-stage control flow. Every routing request pays for an initial SLM pass,
and 34 requests also pay for a second LLM pass. Within the routing run, the
66 SLM-only samples averaged 16.26 seconds, while escalated samples averaged
28.57 seconds.

The escalation breakdown also helps explain the result. Out of the 34 escalated
samples, 23 were escalated because confidence fell below the 0.8 threshold and
11 were escalated because the draft answer could not be parsed reliably. In
other words, part of the latency penalty comes from router conservatism and
output-format failures, not only from difficult questions.

\begin{figure}[htbp]
\centering
\small
\renewcommand{\arraystretch}{1.15}
\begin{tabular}{@{}p{3.3cm}l l r@{}}
\textbf{Metric} & \textbf{System} & \textbf{Relative bar} & \textbf{Value} \\
\hline
Average latency & Routing & \metricbar{routingcolor}{4.00cm} & 20.45 s \\
 & Monolithic & \metricbar{monolithiccolor}{1.82cm} & 9.29 s \\
\\[-0.1em]
Latency p95 & Routing & \metricbar{routingcolor}{4.00cm} & 42.18 s \\
 & Monolithic & \metricbar{monolithiccolor}{1.67cm} & 17.66 s \\
\\[-0.1em]
Total cost & Routing & \metricbar{routingcolor}{3.90cm} & \$0.5913 \\
 & Monolithic & \metricbar{monolithiccolor}{4.00cm} & \$0.6066 \\
\\[-0.1em]
Total energy & Routing & \metricbar{routingcolor}{3.22cm} & 0.0665 kWh \\
 & Monolithic & \metricbar{monolithiccolor}{4.00cm} & 0.0826 kWh \\
\\[-0.1em]
Total CO$_2$ & Routing & \metricbar{routingcolor}{3.22cm} & 25.28 g \\
 & Monolithic & \metricbar{monolithiccolor}{4.00cm} & 31.39 g \\
\\[-0.1em]
LLM call ratio & Routing & \metricbar{routingcolor}{1.36cm} & 34\% \\
 & Monolithic & \metricbar{monolithiccolor}{4.00cm} & 100\% \\
\end{tabular}
\caption{Efficiency and resource comparison on MMLU. Lower is better in all panels.}
\label{fig:mmlu-efficiency}
\end{figure}

\section{Discussion of the Current Result}

These two experiments show that architecture choice changes the operating point
of the system rather than producing a universal winner. If the main objective is
maximum MMLU accuracy, the monolithic \texttt{gemma4-31b} baseline is stronger.
If the objective is to reduce heavy-model usage, energy, and carbon while
keeping accuracy within a usable range, the routing strategy is more attractive.

At the same time, these runs should still be interpreted as a focused pilot
comparison rather than as the final thesis result set. They demonstrate that the
reporting stack, experiment runner, and cost-energy instrumentation can support
meaningful architecture analysis, but they do not yet replace the broader coding
benchmark that the thesis is designed around.

\section{Implemented and Usable Today}

The following pieces are already implemented and useful for development, smoke
testing, and early-stage experimentation:

\begin{itemize}
    \item the local control plane consisting of backend, frontend, MLflow integration, experiment runner, and result reporting;
    \item the canonical model alias system that maps architecture configurations to remote vLLM endpoints;
    \item baseline and hybrid architecture surfaces such as monolithic inference, routing, multi-agent debate, and SLM voting ensembles;
    \item proxy-based and response-level metric propagation for latency, token counts, cost, energy, and emissions; and
    \item web surfaces for dashboarding, experiment launch, host visibility, result comparison, and playground-style prompt testing.
\end{itemize}

These capabilities are sufficient to support platform demonstrations,
architecture debugging, early pilot runs, and documentation of the end-to-end
workflow. They are not yet sufficient to support the final experimental claims
of the thesis.

\section{Preliminary Components Not Yet Valid for Thesis Claims}

Several core ideas of the thesis are still under construction or still require
methodological cleanup:

\begin{itemize}
    \item The coding benchmark is not yet finalized. The intended \texttt{custom\_stratified} benchmark should contain easy, medium, and hard coding tasks, but the current repository loader is still a legacy mixture and therefore cannot be treated as a valid coding benchmark.
    \item The project-specific \texttt{HumanEval} surface is conceptually defined as a human preference workflow, but the current repository file still contains a legacy prototype and is not the final benchmark implementation.
    \item The local end-to-end runner path still exhibits environment and benchmark-loader friction. In practice, this means that a repository-level ``run everything'' path is not yet stable enough to be cited as a finished validation pipeline.
    \item The project-proposed swarm family --- decentralized blackboard, entropy blackboard, and pure swarm --- is implemented as a promising architecture direction, but not yet supported by a finalized benchmark and valid comparative evidence.
\end{itemize}

For these reasons, the current experimental outputs are treated as exploratory
outside the specific MMLU comparison reported above. No broader quantitative
claim is made in this term report for the unfinished benchmark surfaces. This is
intentional: fabricated or unstable numbers would weaken the thesis rather than
strengthen it.

\section{Planned Coding Benchmark Protocol}

The main benchmark contribution planned for the thesis is a coding-focused
dataset organized by difficulty. The benchmark will be centered on the
\texttt{custom\_stratified} track and will compare architectures on practical
coding tasks rather than on a single public multiple-choice benchmark. The
target structure is shown in Table \ref{tab:coding-benchmark-plan}.

\begin{table}[htbp]
    \centering
    \caption{Planned size targets for the coding benchmark.}
    \label{tab:coding-benchmark-plan}
    \small
    \begin{tabularx}{\linewidth}{@{} l c c c X @{}}
        \toprule
        \textbf{Phase} & \textbf{Easy} & \textbf{Medium} & \textbf{Hard} & \textbf{Purpose} \\
        \midrule
        Pilot set & 10 & 10 & 10 & Quick validation of scoring, prompts, and end-to-end runner behavior. \\
        Main comparison set & 50 & 50 & 50 & Minimum thesis-grade automated comparison target across architectures and model tiers. \\
        Stronger final set & 100 & 100 & 100 & Ideal stretch target if authoring, checking, and hidden-test preparation are completed in time. \\
        \bottomrule
    \end{tabularx}
\end{table}

Each task record should contain at least an identifier, a difficulty label, a
prompt, optional starter code, public and hidden tests, a topic label, and
source metadata. The benchmark should support evaluation by exact output or
unit tests depending on the task. The main comparison axes will be:
architecture, SLM choice, LLM choice, escalation behavior, total cost, energy
consumption, carbon emissions, and EATS. This structure makes the benchmark
suitable for the thesis focus on coding skill rather than on general-purpose
language understanding alone.

\section{Thesis-Grade Evaluation Protocol}

Once the benchmark is finalized, the comparative protocol will proceed in three layers:

\begin{itemize}
    \item \textbf{Monolithic baseline:} a single stronger model answers every query and provides the upper baseline for quality and the likely upper bound for cost.
    \item \textbf{Hybrid and ensemble comparison:} routing, multi-agent, and ensemble systems are compared against the monolithic baseline under the same coding tasks and the same reporting stack.
    \item \textbf{Project-proposed architecture study:} the swarm family is evaluated after the benchmark and runner paths are stable enough to support valid comparison.
\end{itemize}

All comparisons will record accuracy or task success, latency, input/output tokens, total cost, API/infrastructure cost split, LLM-call ratio, energy, emissions, and EATS. Human preference evaluation will be added in a later phase through the UI-backed benchmark surface and should be treated as a complementary study rather than as a replacement for the coding benchmark.

\section{Pending Result Artifacts}

The following quantitative artifacts are intentionally left as placeholders until stable runs are completed:

\begin{table}[htbp]
    \centering
    \caption{Pending result artifacts for later insertion.}
    \label{tab:pending-results}
    \small
    \begin{tabularx}{\linewidth}{@{} l X @{}}
        \toprule
        \textbf{Artifact} & \textbf{Current status} \\
        \midrule
        Accuracy / latency / cost summary table & Pending insertion after valid automated benchmark runs. \\
        Energy / CO$_2$ / EATS comparison table & Pending insertion after stable proxy-backed measurements are collected. \\
        Architecture comparison figure & Pending insertion after enough runs exist for a meaningful multi-architecture plot. \\
        Human preference result table & Pending insertion after the UI-backed preference workflow is finalized and used. \\
        \bottomrule
    \end{tabularx}
\end{table}

\begin{figure}[htbp]
    \centering
    \fbox{\parbox{0.85\linewidth}{\centering Pending figure placeholder: insert a final architecture comparison plot after the coding benchmark, measurement stack, and experiment pipeline are all stable enough for valid thesis claims.}}
    \caption{Reserved placeholder for the final comparative experiment figure.}
    \label{fig:pending-comparison}
\end{figure}
